\section{Related Work}
\label{sec:related}

\subsection{Automated Approximate-Accelerator Design}

Approximate computing has produced a broad range of circuit-, architecture-, and system-level techniques for trading output accuracy against hardware efficiency \cite{xu2016approxsurvey}. Recent surveys of approximate-computing DSE further show that automated exploration commonly combines arithmetic approximation with evolutionary, learning-based, or search-based optimization \cite{saeedi2024dse,silva2024mapping}. Approximate arithmetic libraries such as EvoApproxLib and related collections provide large sets of implementations with different error and hardware characteristics \cite{mrazek2017evoapprox,hanif2017quad,ullah2018smapproxlib}.

Automation has been studied at several abstraction levels. Approximate logic synthesis directly transforms Boolean or gate-level implementations under explicit error constraints \cite{miao2013als,venkataramani2020logicsynthesis}. Axilog provides language-level support for identifying relaxable portions of RTL hardware \cite{yazdanbakhsh2015axilog}, while Chisel selects approximate operations according to application-level reliability and accuracy constraints \cite{misailovic2014chisel}. These studies establish approximation placement as an important design decision in automated approximate computing.

At the accelerator level, ABACUS, AutoAx, and AxHLS automate the construction of larger approximate systems from approximate operations or components \cite{nepal2014abacus,mrazek2019autoax,castrogodinez2020axhls}. FPGA-oriented and multi-stage frameworks further extend automated approximate design across different implementation flows \cite{prabakaran2020approxfpgas,prabakaran2023xelfpgas,meng2023hedals,meng2025appresub}. Search-based approaches have also been introduced to navigate large approximate design spaces, including Monte-Carlo-tree-search-based accelerator synthesis and learning-assisted approximate-circuit DSE \cite{awais2021mcts,awais2021ldax}.

A common accelerator-level formulation assigns one library implementation to each operation exposed to DSE. AutoAx represents an accelerator configuration through approximate-component assignments \cite{mrazek2019autoax}, and ApproxPilot similarly optimizes arithmetic-unit choices across accelerator operations \cite{zhang2025approxpilot}. For large RTL targets, approximation placement and unit assignment become closely coupled because the most effective design may retain Exact implementations at a substantial subset of legally configurable sites. This motivates a search space in which Exact and approximate implementations are treated uniformly as candidate states at each configurable arithmetic operation.

\subsection{Learning-Based PPA/QoR Modeling and Transfer}

Machine learning has become an established means of reducing repeated EDA evaluations, spanning performance prediction, physical-design estimation, black-box optimization, and learned circuit representations \cite{huang2021mleda}. GNN-based EDA further exploits the natural graph structure of RTL, netlists, and circuit data \cite{sanchez2023gnneda}.

For high-level and accelerator design, graph models have been applied to latency and performance estimation and to learning-assisted design-space exploration \cite{ustun2020delay,wu2022hls_gnn,wu2023ironmanpro}. Related models address power and PPA prediction \cite{lin2022powergear,lin2023hlpow,zhao2025pdnnet}, while structural circuit representation learning provides reusable graph embeddings for downstream EDA tasks \cite{xu2022sns,li2022deepgate,shi2023deepgate2}.

Learning-based modeling has also become central to approximate-accelerator exploration. AutoAx uses conventional regression models to estimate accelerator-level hardware cost and output quality \cite{mrazek2019autoax}. ApproxPilot introduces graph-based models together with critical-path-aware delay prediction \cite{zhang2025approxpilot}. These methods construct target-specific prediction models from sampled accelerator configurations.

ApproxGNN develops a complementary transfer-oriented approach. It learns approximate-component embeddings from a reusable corpus of 15.9K configurations generated from 159 random image-convolution kernels and evaluates transfer to unseen Small and Large Gaussian filters \cite{vlcek2025approxgnn}. This demonstrates that learned component representations can be reused across related accelerator structures and reduces the need to train every predictor entirely from scratch.

Unit characterization provides another reusable source of information. Area, Power, Delay, and error properties of arithmetic units can be obtained independently of a particular accelerator and incorporated into its graph representation. The resulting system model can use target-system labels to capture composition, synthesis, and topology-dependent interactions while retaining reusable local hardware information.

\subsection{Sequential Optimization with Expensive Evaluations}

Design-space exploration under expensive evaluation budgets has been widely studied using surrogate and sequential optimization. Bayesian optimization provides a general framework for choosing new evaluations from the current predictive model \cite{shahriari2016bo,knowles2006parego,wang2023bo}, while trust-region and tree-based approaches extend model-based optimization to larger and more structured spaces \cite{eriksson2019turbo,wang2020lamcts,zhao2021mopartitions}.

Sample-efficient optimization has also been applied directly to hardware design. HyperMapper supports multiobjective design-space exploration with expensive and potentially infeasible evaluations and has been demonstrated on hardware accelerator tuning \cite{nardi2019hypermapper}. BOOM-Explorer combines active sampling, learned surrogate modeling, and multiobjective Bayesian optimization for processor microarchitecture DSE \cite{bai2021boomexplorer}. Accelerator-rich systems have similarly been explored with active-learning-based refinement around promising Pareto regions in ARS-Flow and its extended formulation ARS-Flow 2.0 \cite{huang2024arsflow,huang2025arsflow2}.

Related ideas appear in accelerator design, approximate arithmetic, circuit optimization, and Pareto-front identification. Bayesian optimization has been applied to accelerator DSE \cite{reagen2017bayesopt}, while active or surrogate-guided methods have been investigated for approximate adders and circuit design \cite{ma2019adder_active,ma2019cadbo}. Other work considers surrogate-assisted multiobjective search, Pareto identification, and active evaluation in hardware and circuit optimization \cite{geng2022gnp,belwal2022qpir,cao2025roseopt,chen2026yield,yin2022activeyield}.

These methods establish the value of allocating expensive evaluations according to an evolving optimization state. Approximate-accelerator DSE adds a distinctive coupling between this allocation problem and the underlying configuration space: the optimizer must simultaneously determine approximation placement, arithmetic-unit binding, and which target-system configurations should receive the next synthesis evaluations. ApproxPilot-SE integrates these decisions within a single closed-loop DSE process under a fixed target-specific evaluation budget.